\documentclass[11pt]{article}
\usepackage[a4paper,margin=1in]{geometry}
\usepackage{fontspec}
\usepackage[boldfont=false]{xeCJK}
\setCJKmainfont{FandolSong-Regular.otf}
\usepackage{amsmath}
\usepackage{amssymb}
\usepackage{array}
\usepackage{booktabs}
\usepackage{graphicx}
\usepackage{adjustbox}
\usepackage{enumitem}
\setlist[itemize]{topsep=2pt,partopsep=0pt,parsep=0pt,itemsep=2pt,leftmargin=1.4em}
\setlist[enumerate]{topsep=2pt,partopsep=0pt,parsep=0pt,itemsep=2pt,leftmargin=1.6em}
\usepackage{parskip}
\usepackage{fvextra}
\fvset{breaklines=true,breakanywhere=true,fontsize=\small}
\setlength{\parindent}{0pt}

\begin{document}

\begin{center}
  {\LARGE\bfseries Immunity}\\[0.35em]
  {\large A Level Biology}
\end{center}
\vspace{0.6em}

\section*{Phagocytes — the first defence}

Your \textbf{immune system} 免疫系统 is the set of cells that defends the body against \textbf{pathogens} 病原体. The first cells to act are \textbf{phagocytes} 吞噬细胞, a group of white blood cells that includes \textbf{macrophages} 巨噬细胞 and \textbf{neutrophils} 中性粒细胞.

A phagocyte destroys a pathogen by \textbf{phagocytosis} 吞噬作用: it surrounds the pathogen, takes it inside in a vesicle, and digests it with enzymes. After this, a macrophage displays parts of the pathogen on its own surface, ready to alert other immune cells.

\section*{Antigens: self and non-self}

An \textbf{antigen} 抗原 is a molecule (usually a protein) on a cell surface that the immune system can recognise.

\begin{itemize}
\item \textbf{self antigens} 自身抗原 are the body's own markers. The immune system learns to ignore them.

\item \textbf{non-self antigens} 非自身抗原 are foreign, for example the antigens on a pathogen. These trigger an immune response.

\end{itemize}

\section*{The primary immune response}

The first time a new pathogen enters, the body makes a slow \textbf{immune response} 免疫反应. The main steps are:

\begin{enumerate}
\item a macrophage engulfs the pathogen and displays its antigen.

\item \textbf{T-helper cells} 辅助性T细胞 recognise that antigen and become active. They release chemicals that switch on other cells.

\item \textbf{B-lymphocytes} 淋巴细胞 with a matching shape are selected. They divide to form \textbf{plasma cells} 浆细胞, which pour out \textbf{antibodies} 抗体, and \textbf{memory cells} 记忆细胞.

\item \textbf{T-killer cells} 杀伤性T细胞 destroy the body's own cells that have been infected.

\end{enumerate}

\section*{Memory cells and long-term immunity}

Memory cells stay in the body for years after the infection is over. If the \textbf{same} pathogen enters again, the memory cells start a \textbf{secondary} immune response that is much faster and larger than the first. The pathogen is destroyed before it can make you ill. This is what we mean by long-term immunity.

\section*{Antibodies}

An antibody is a Y-shaped protein made by plasma cells. The two tips of the Y are \textbf{antigen-binding sites} 抗原结合位点. Each site has a special shape, called the \textbf{variable region} 可变区, that fits one antigen only, like a lock and key.

Antibodies help in several ways: they stick to antigens, clump pathogens together so they are easier to deal with, mark pathogens so phagocytes find them, and block harmful toxins.

\section*{Monoclonal antibodies}

A \textbf{monoclonal antibody} 单克隆抗体 is a single type of antibody, all identical. They are made by the \textbf{hybridoma} 杂交瘤 method:

\begin{enumerate}
\item an animal is given an antigen, so it makes B-lymphocytes that produce the wanted antibody.

\item these B-lymphocytes are fused with tumour cells, which divide endlessly.

\item the fused cell (the hybridoma) both makes the antibody \textbf{and} divides without stopping, producing large amounts of one identical antibody.

\end{enumerate}

Monoclonal antibodies are used in the \textbf{diagnosis} 诊断 of disease (to detect a specific molecule, as in a pregnancy test) and in treatment (to carry drugs to specific target cells, such as cancer cells).

\section*{Types of immunity}

Immunity can be active or passive, and natural or artificial.

\begin{itemize}
\item \textbf{active immunity} 主动免疫 — your own body meets an antigen and makes its own antibodies and memory cells. It is slow to start but long-lasting.

\item \textbf{passive immunity} 被动免疫 — ready-made antibodies are given to you from outside. It works at once but does not last, because there are no memory cells.

\item \textbf{natural immunity} 天然免疫 — gained in a natural way (active: after an infection; passive: antibodies passed from mother to baby).

\item \textbf{artificial immunity} 人工免疫 — gained on purpose (active: your body is made to respond to a safe dose of antigen; passive: you are injected with ready-made antibodies).

\end{itemize}

\section*{Vaccination}

A \textbf{vaccine} 疫苗 contains antigens — often a dead or weakened pathogen, or part of one. The antigens trigger a primary immune response and make memory cells, so you gain long-term immunity \textbf{without} becoming ill.

\textbf{Vaccination} 疫苗接种 programmes can control the spread of a disease across a population. If enough people are vaccinated, the pathogen cannot pass easily from person to person. This protects even the people who are not vaccinated, an effect called \textbf{herd immunity} 群体免疫.


\end{document}
